\documentclass{jasb}
\usepackage[utf8]{inputenc}
\usepackage[T1]{fontenc}
\usepackage{geometry}
\usepackage{graphicx}
\usepackage{float}
\usepackage{hyperref}
\usepackage[english,indonesian]{babel}

\title{Optimasi Pemrosesan Template LaTeX Menggunakan Pendekatan Berbasis AI}
\author{Academic Support Agent}
\date{\today}

\begin{document}
\maketitle

\begin{abstract}

\end{abstract}

\chapter{PENDAHULUAN}


\chapter{TINJAUAN PUSTAKA}


% Original content (commented out):
% \documentclass[english]{jasb}
% %*************************************************************************
% % Nastavení čísel 
% \def \year {2021}
% \def \volume {1}
% \def \issue {2}
% \def \firstpagenumber {00}
% \def \lastpagenumber {04}
% \setDoi{aaa}
% %*************************************************************************
% IT IS RECOMMENDEDTO HAVE FIGURES IN A SEPARATE DIRECTORY "figures/"
\graphicspath{ {figures/} }
%*************************************************************************
\title{Template for Submission of Manuscripts to Journal of Alumni Scientiae Bohemicae Society}
\date{}

\names{
    \name[*]{Grey, Gandalf}{1}, \name{Potter, Harry}{2}
}
\institutions{
    \institution{1}{Fellowship of the Ring, Middle-earth}
    \institution{2}{Hogwarts School of Witchcraft and Wizardry, Scotland}
}
\correspondingEmail{gandalf.grey@fellowship.me}
\abstractt{
    This template serves as a guide to prepare manuscripts for submission. 
    Authors are encouraged to read the instructions below to avoid unnecessary 
    corrections later.
}
\keywords{typesetting, article, xelatex, bibtex}

%*******************************************************************************
% Author's packages and macros
%*******************************************************************************


%*******************************************************************************
\begin{document}
   \maketitle 
   %*******************************************************************************
% DELETE FROM HERE
%*******************************************************************************
   This template serves as a guide to prepare manuscripts for submission. The entire manuscript (text, tables, graphics, bibliography) should be submitted as one zip file. Inserting graphics and tables close to the point at which they are discussed in the text of the manuscript will benefit both the reviewers and the readers. It also possible to send the graphics and tables in separate files, if it is more convenient for the author.

This template is not trying to produce the final formatting that will be applied only after acceptance during the editorial process. The predefined \LaTeX{} commands and preloaded packages are there not only to help you, the author, but they also aid the editor and reviewer in the orientation in the bulk of the text, thus their use is encouraged. To see the final formatting please check the \href{http://journal.asbsociety.com/articles/}{article example on the journal page}.
\section*{Important Manuscript Requirements}

\begin{enumerate}
\item All authors must be listed with current contact information and a direct/valid e-mail address for each.
\item Articles must contain an Abstract and optionally a Table of Contents (TOC) graphic. The TOC graphic is shown next to the abstract.
\item Articles that exceed the norm of 5-8 Journal pages (roughly 4500-8000 words, references not included), often contain a level of detail that is not appropriate for a broad audience. Please consider moving of some parts into the Supporting Information.
\item Smaller, ``communication-like'' articles (articles not exceeding 2200 words) are encouraged.
\item All sections must be properly labeled, please adhere to the template.
\item If you want to submit a Supporting Information file (SI), please use the PDF format and number the pages consecutively. For data, CSV can be used. For groups of these files, ZIP files may be used.
\item Appendices belong into the SI.
\item Use of excessive self-citations is not permitted. (i.e.\ citations of authors of the manuscript)
\item Disclaimers are not allowed in manuscripts or in the SI.
\item Trademarked items are not permitted.
\item Dedications may appear only in the ``ACKNOWLEDGEMENTS'' section.
\end{enumerate}
Instructions for using the template follow on the next page.
\newpage

\section*{Using the Template}
This template is intended to be used directly to ease the writing process for the first-time author. Simply follow the guidelines below and rewrite the placeholder text. All the sections are obligatory, please do not skip any.
\begin{enumerate}
\item The article is to be compiled using XeLaTeX with BibLaTeX (\texttt{biber} backend) and packages present in the TeX Live distribution. Please make sure to configure your computer accordingly.
\item When typing your manuscript directly into the template, delete all contents of the \texttt{submission\_template.tex} file between \textsc{delete from here} and \textsc{delete until here}.
\item The \texttt{natsci} package will be used to typeset the final article and should be enclosed with this template. It loads most packages that can authors are expected to use.
\item You are strongly encouraged to use the \texttt{natsci} package style guide (attached) when typesetting your submission.
\item Write in as plain \LaTeX{} as possible. Use the predefined  commands for sections.
\item For additional functionality, first check the packages loaded by the \texttt{natsci} package and only if these are insufficient, add your own in the designated part fo the header. Make sure that these do not collide with any of the packages loaded by \texttt{natsci}.
\item Include all your personal macros in the designated part fo the header. Do not add your own \texttt{.sty} or \texttt{.cls} files.
\item Insert graphics using the \texttt{figure} environment.
\begin{enumerate}
    \item Please insert graphics in the exact size in which they are to be included in the final article. Please ensure that it is perfectly legible and crisp (maximum width is 17 cm). A convenient option might be setting parameter 
        \texttt{width=\\linewidth} (1-column figure) or \texttt{width=\\textwidth} (page-wide figure).
    \item NOTE: If you are submitting a Table of Contents graphic, please insert the graphic at the end of the file.
    \item Vector graphics (recommended)
    \begin{enumerate}
        \item Format: PDF, EPS, drawn using the \texttt{tikz} package
    \end{enumerate}
    \item Raster graphics
    \begin{enumerate}
        \item Resolution: 300 dpi for color, 600 dpi for grayscale art, and at 1200 dpi for black and white line art (e.g. graphs). Ideally, include graphics with \texttt{scale=1} to avoid errors.
        \item Format: PNG, JPG
    \end{enumerate}
\end{enumerate}
\item Adhere to the format of references described at the end of the template below.
\item Delete these instructions and manuscript requirements above.
\item Make sure to upload your \texttt{.tex}, \texttt{.bib} and all figure files.
\item Proof the manuscript to ensure that all parts of the manuscript are present and clearly legible.
\end{enumerate}
Template follows on the next page.
\newpage
%*******************************************************************************
% DELETE UNTIL HERE
%*******************************************************************************
\textbf{TITLE.} The title should accurately, clearly, and concisely reflect the emphasis and content of the paper. The title must be brief and grammatically correct. The title should NOT contain abbreviations. It is also not advisable to use words such as ``first'', ``novel'' nor the title should contain any part or series number.

\textbf{AUTHOR NAMES.} Include in the by-line all those who have made substantial contributions to the work, even if the paper was actually written by only one person. Use first names, initials, and surnames (e.g., Karel J. Erben) or first initials, second names, and surnames (e.g., K. Jaromír Erben). Do not use only initials with surnames (e.g., K. J. Erben) because this causes indexing and retrieval difficulties and interferes with unique identification of an author. Do not include professional or official titles or academic degrees. At least one author must be designated with an asterisk* as the author to whom correspondence should be addressed.

\textbf{AUTHOR ADDRESS.} The affiliation should be the institution where the work was conducted. If the present address of an author differs from that at which the work was done, indicate with a number$^1$ and give the Present Address under Author Information.  If more than one address, use numbers to match author names to address(es).

\textbf{KEYWORDS} Provide significant keywords to aid the reader in literature retrieval.

\begin{abstract}
    All manuscripts must be accompanied by an abstract. The abstract should briefly state the problem or purpose of the research, indicate the theoretical or experimental plan used, summarize the principal findings, and point out the major conclusions. Abstract length is one paragraph. Be brief and clear.
\end{abstract}

\section{Heading}
TEXT. This is the place for the text of the article. The article should start with an introductory part (``INTRODUCTION'') explaining the problem at hand. It is eagerly recommended to add such part as the journal is aimed towards broader scientific audience. This then may be followed by a part discussing theory, experiments, or methodology (``THEORETICAL PART'', ``EXPERIMENTAL PART'', ``METHODOLOGY''), whichever suits your field of study the best. It is a good practice for the article to culminate in two sections with formal list of results (``RESULTS'') and a section discussing the said results (``DISCUSSION''). These two may be joined (``RESULTS AND DISCUSSION''), but especially in natural sciences it often leads to obfuscation of the message the article tries to deliver, thus it is not recommended. The last part is a brief overview of the results and possibly the outlook of future research (``CONCLUSION''). Be wary of the fact that the ``CONCLUSION'' part and ``ABSTRACT'' parts are different not only in length, but also in contents. For more inspiration, check the article example on the journal page.

\section{Displayed equations}
Displayed equations can be inserted where desired using the \texttt{equation environment}. Equations are numbered with Arabic numerals in parentheses on the right-hand side as automatically done by LaTeX using the \texttt{equation} environment. For lengthy equations, groups of equations, arrays of equations and other special cases, you can use the environments from the \texttt{amsmath} package like \texttt{gather}, \texttt{align} and \texttt{multline}. You can use \LaTeX{} labels to refer to equations, like this~\eqref{eq:xdefintion}.
\begin{equation}
    x=y
    \label{eq:xdefintion}
\end{equation}
Inline equations can be included using the ``single dollar'' $x=y$ or ``escaped parentheses'' \(x=y\) notation. If unnumbered display equations are required, the \texttt{equation*} environment and \texttt{\textbackslash[x=y\textbackslash]} should be used. Please avoid the old \TeX ``double dollar'' notation \texttt{\$\$x=y\$\$}.

Long equations can be typesetted to the whole width of the page using the 
\texttt{strip} environment, if needed on multiple lines as 
in example~\eqref{eq:wide_equation}.
 \begin{strip}
   \begin{equation}
   \begin{split}
      OutcomeVariable_{it} = \alpha &+ \sum_{q = 1,2} \beta_1^q (StimulusPayments) \cdot 1 \left[ReceivedPayments=q\right] + \\
      &+ \beta_2^1(WeekAfterSimulusPayments)_t \cdot 1 \left[ReceivedPayments=q\right] + \epsilon\,,
       \label{eq:wide_equation}
   \end{split}
   \end{equation}
   \end{strip}

\section{Figures}
Figures must be inserted using the \texttt{figure} (or \texttt{figure*} for 2-column) environment. Each figure must have a caption that includes a brief description. All figures must be mentioned in the text consecutively (like Fig.~\ref{fig:myfigure}). The caption should be understandable without reference to the text. Whenever possible, place the key to symbols in the artwork, not in the caption. To insert the graphic into the figure, be sure it is already sized appropriately and place it before the figure caption.
% \begin{figure} [H]
%     \centering
%     \includegraphics[width=\linewidth]{qheom_contour.pdf}
%     \caption{One column figure caption.
%     }
%     \label{fig:myfigure}
% \end{figure}

% \newpage

\begin{figure*}[ht]
    \centering
    \includegraphics[width=0.9\textwidth]{qheom_contour.pdf}
    \caption{Two column figure caption.
    }
    \label{fig:myfigure}
\end{figure*}

\section{Tables}
Each table must have a brief (one phrase or sentence) title that describes its contents. The title should be understandable without reference to the text. Put details in footnotes below the table, not in the title. Like figures, tables should be mentioned in the text consecutively (like Table~\ref{tab:mytable}). 
Use \texttt{table} or \texttt{table*} environment for 1- or 2-column tables 
respectively.

\begin{table}[H]
    \centering
    \caption{Brief self-explanatory title of a one-column wide table.}
    \begin{tabular}{ p{3cm} p{3cm} }
    % \begin{tabular}{ |p{3cm}||p{3cm}|}
        % \hline
        \toprule
        \multicolumn{2}{c}{Country List} \\
        % \multicolumn{2}{|c|}{Country List} \\
        % \hline
        Country Name     or Area Name& ISO ALPHA 2 Code \\
        \midrule
        Afghanistan   & AF\\
        Aland Islands&   AX  \\
        Albania &AL \\
        Algeria    &DZ\\
        American Samoa&   AS \\
        Andorra& AD \\
        Angola& AO \\
        % \hline
        \bottomrule
    \end{tabular}
    % Delete from here if footnotes unnecessary
    \footnotesize

    \vspace{0.5em}
    Footnotes below the table giving a more detailed explanation
    % Delete to here if footnotes unnecessary
    \label{tab:mytable_narrow}
\end{table}

\begin{table*}[ht]
    \centering
    \caption{Brief self-explanatory title of a two-column wide table.}
    \begin{tabular}{ p{3cm} p{3cm} p{3cm} p{3cm}  }
    % \begin{tabular}{ |p{3cm}||p{3cm}|p{3cm}|p{3cm}|  }
        \toprule
        % \hline
        \multicolumn{4}{c}{Country List} \\
        % \multicolumn{4}{|c|}{Country List} \\
        % \hline
        Country Name     or Area Name& ISO ALPHA 2 Code &ISO ALPHA 3 Code&ISO numeric Code\\
        \midrule
        % \hline
        Afghanistan   & AF    &AFG&   004\\
        Aland Islands&   AX  & ALA   &248\\
        Albania &AL & ALB&  008\\
        Algeria    &DZ & DZA&  012\\
        American Samoa&   AS  & ASM&016\\
        Andorra& AD  & AND   &020\\
        Angola& AO  & AGO&024\\
        % \hline
        \bottomrule
    \end{tabular}
    % Delete from here if footnotes unnecessary
    \footnotesize

    \vspace{0.5em}
    Footnotes below the table giving a more detailed explanation
    % Delete to here if footnotes unnecessary
    \label{tab:mytable}
\end{table*}


\section*{Associated content}
Supporting Information
The Supporting Information is available on the journal website.
brief description (file type, i.e., PDF or CSV)
brief description (file type, i.e., PDF or CSV)

\section*{Acknowledgements}
This is generally the place for dedications, to acknowledge people, or places. You may state grant numbers and sponsors here as well.
\begin{figure} [H]
    \centering
    \includegraphics[width=\linewidth]{qheom_contour.pdf}
    \caption{One column figure caption.
    }
    \label{fig:myfigure}
\end{figure}

\section*{Abbreviations}
Write down all the abbreviations used to help the reviewer (e.g.: NMR = Nuclear Magnetic Resonance, BASIC = Beginner's All-Purpose Symbolic Instruction Code, RADAR = Radio Detection and Ranging, SARS = Severe Acute Respiratory Syndrome).

\dpic{qheom_contour.pdf}{qheom_contour.pdf}{Two figures side-by-side. Use only if you have a good reason for that.}{fig:side-by-side-qheom}

\section*{Citing  references}
For citations, the used package is \texttt{biblatex} with \texttt{biber} backend. The citations are then imported from a separate \texttt{.bib} file, which can be automatically generated from a citation manager, like Zotero, Mendeley or EndNote.

To add a reference, please use the superscript notation, e.g.:

This formula is called after Enrico Fermi, who popularized the research of Paul Dirac that was the one actually obtaining it when trying to describe the emission and absorption of radiation using quantum theory during the year 1927.\supercite{Dirac1927,Pauli}

The reference itself belongs here, below ``REFERENCES'' heading:


Note: The first reference is pointing to an article by Paul Adrien Maurice Dirac called ``The quantum theory of the emission and absorption of radiation.'' in a journal named Proceedings of the Royal Society of London, Series A (abbreviated by Standard Abbreviation (ISO4): Proc. R. Soc. Lond.). It was published in 1927 in journal's 114th volume on the page 243 and has a Digital Object Identifier (DOI) 10.1098/rspa.1927.0039. It is heavily encouraged to use DOI's if the reference has it as it greatly helps the reader to find said reference by using doi.org/''digital object identifier'' webpage.

The second reference is pointing to a book part written by Wolfgang Pauli in a book named Pauli Lectures on Physics, Vol. 5, in which the editor was Charles Paul Enz. It was published by Dover Publications, Inc. in Mineola, state New York, in the year 1973, and the book part itself is situated through the pages 147 and 151. The last information given is the ISBN-13 code.


%*******************************************************************************
% DELETE UNTIL HERE
%*******************************************************************************
\printbibliography
\end{document}


\end{document}